% Required package: \usepackage{booktabs}
\subsection{Motivation for Systolic-Array Power Management}
\label{sec:component-selection}

To identify a suitable power-management target, we augmented ONNXim with
cycle counters for the systolic array (SA), the core memory-request FIFO, the
on-chip network (NoC), and DRAM. The counters are observational: they do not
modify scheduling, instruction issue, or simulated execution timing. SA
idleness denotes cycles in which the array executes no SA operation. NoC
idleness denotes cycles with no demand packet in flight, while DRAM idleness
denotes cycles in which all channels have no outstanding demand request.
Therefore, the NoC and DRAM values represent demand occupancy rather than
physical link, bank, or command-bus utilization.

Table~\ref{tab:component-characterization} reports the measured idle-cycle
fractions for each workload. The SA is idle for 44.64--75.36\% of execution
across the evaluated suite. The opportunity persists across both CNN and
transformer workloads rather than arising from a single model. In particular,
the larger transformer configurations exhibit approximately 73--75\% SA
idleness even though their NoC and DRAM demand-idle fractions are only
6--15\%. Thus, the array frequently waits while the surrounding memory system
remains occupied, motivating local SA power management rather than
coarse-grained shutdown of the complete NPU or memory subsystem.

\begin{table*}[t]
  \centering
  \caption{Measured idle-cycle fractions across the workload suite. Larger
  OPT and LLaMA configurations use representative decoder-layer execution;
  OPT-125M is evaluated end-to-end.}
  \label{tab:component-characterization}
  \begin{tabular}{lrrrr}
    \toprule
    Workload & SA idle (\%) & Request FIFO empty (\%) &
    NoC demand-idle (\%) & DRAM demand-idle (\%) \\
    \midrule
    AlexNet            & 75.36 & 62.25 & 54.00 & 54.04 \\
    ResNet-18          & 53.64 & 98.89 & 94.49 & 94.65 \\
    ResNet-34          & 60.68 & 98.84 & 93.94 & 94.12 \\
    ResNet-50          & 63.63 & 95.54 & 81.59 & 82.04 \\
    ResNet-101         & 68.18 & 95.83 & 80.24 & 80.86 \\
    VGG-11             & 65.51 & 53.14 & 37.41 & 37.46 \\
    VGG-16             & 63.78 & 59.70 & 48.43 & 48.48 \\
    VGG-19             & 63.72 & 64.47 & 52.79 & 52.85 \\
    Wide-ResNet-50-2   & 68.10 & 96.14 & 85.46 & 85.81 \\
    GPT-2              & 44.64 & 62.27 & 34.05 & 34.21 \\
    OPT-125M (full)    & 73.27 & 65.43 & 25.43 & 26.58 \\
    OPT-1.3B           & 75.05 & 25.44 & 14.47 & 14.76 \\
    OPT-2.7B           & 74.56 & 22.14 & 14.46 & 14.70 \\
    OPT-6.7B           & 74.75 & 22.84 & 10.26 & 10.36 \\
    OPT-66B            & 73.59 & 26.11 &  6.12 &  6.15 \\
    LLaMA-3-8B         & 74.16 & 23.06 & 10.48 & 10.57 \\
    \midrule
    \textbf{Cycle-weighted aggregate}
                       & \textbf{72.32} & \textbf{31.04} &
                         \textbf{12.56} & \textbf{12.70} \\
    \bottomrule
  \end{tabular}
\end{table*}

Across all runs, the SA is idle for 649,866,804 of 898,576,188 observed
core-cycles, corresponding to a cycle-weighted idle fraction of 72.32\%.
By comparison, the cycle-weighted NoC and DRAM demand-idle fractions are
12.56\% and 12.70\%, respectively. The SA therefore offers a large and
workload-consistent local idle opportunity while remaining active for
27.68\% of observed core-cycles. Based on this combination of frequent
idleness, regular participation in useful matrix computation, and a clear
local control boundary, we select the SA as the target of the proposed
fine-grained power-management mechanism.
